记
\begin{equation*}
M_W\equiv m_Wc^2,\qquad M_Z\equiv m_Zc^2,
\end{equation*}
其中 $M_W$ 与 $M_Z$ 都是能量。$W$ 与 $Z$ 的树级费米子宽度其实是同一个二体衰变母式的两次参数代入。先考虑质量能量为 $M_V$ 的有质量矢量玻色子 $V$ 衰变为两个 Dirac 末态 $f_1,\bar f_2$ 的一般手征顶角
\begin{equation}
 \mathcal M_\lambda
 =\epsilon_\mu^{(\lambda)}(P)\,
 \bar u(p)\gamma^\mu\!
 \left(c_LP_L+c_RP_R\right)v(k),
 \qquad p^2=m_1^2c^2,\quad k^2=m_2^2c^2,
 \label{eq:generic-vector-ff-amplitude-dp91}
\end{equation}
其中 $c_L,c_R$ 是已经吸收规范耦合与内部量子数的无量纲手征系数。允许两个末态质量分别为 $m_1,m_2$ 时，树级宽度在任何质量比下都由精确迹--相空间母式
\begin{equation}
 \boxed{
 \Gamma_{V\to12}
 =\frac{N_c^f}{2(M_V/c^2)\hbar}
 \int\dd\Phi_2\,
 \frac13\left(-\eta_{\mu\nu}+\frac{P_\mu P_\nu}{M_V^2/c^2}\right)
 \Tr\!\left[(\slashed p+m_1c)\Gamma^\mu
 (\slashed k-m_2c)\bar\Gamma^\nu\right],}
 \label{eq:generic-vector-ff-exact-width-dp92}
\end{equation}
给出，其中 $\Gamma^\mu=\gamma^\mu(c_LP_L+c_RP_R)$，$\bar\Gamma^\nu=\gamma^0\Gamma^{\nu\dagger}\gamma^0$。这是当前树级理论层级中的一般结果。为把质量依赖显式化，定义
\begin{equation*}
 r_i\equiv\frac{m_i^2c^4}{M_V^2},\qquad
 \lambda_{12}\equiv\lambda(1,r_1,r_2)
 =(1-r_1-r_2)^2-4r_1r_2,
\end{equation*}
其中 $\lambda$ 是二体运动学的 Källén 函数。完成 Dirac 迹和精确二体相空间积分后得到闭式母式
\begin{equation}
 \boxed{
 \Gamma_{V\to12}
 =N_c^f\frac{M_V}{48\pi\hbar}\sqrt{\lambda_{12}}
 \left\{
 (|c_L|^2+|c_R|^2)
 \left[2-r_1-r_2-(r_1-r_2)^2\right]
 +12\,\operatorname{Re}(c_Lc_R^*)\sqrt{r_1r_2}
 \right\}.}
 \label{eq:generic-vector-ff-exact-closed-dp92}
\end{equation}
这个表达式在树级、两体末态的理论层级内不要求轻质量近似；阈值由 $\lambda_{12}\ge0$ 自动编码。

对 $W,Z$ 的轻费米子部分宽度，控制量
\begin{equation*}
 \epsilon_m\equiv\frac{\max(m_1,m_2)c^2}{M_V}\ll1
\end{equation*}
允许在\eqnrefs{eq:generic-vector-ff-exact-width-dp92} 中取 $m_1,m_2\to0$。此时无偏振结果化为
\begin{equation}
 \boxed{
 \Gamma(V\to f\bar f)
 =N_c^f\frac{M_V}{24\pi\hbar}
 \left(|c_L|^2+|c_R|^2\right)
 +\mathcal O(\epsilon_m^2).}
 \label{eq:generic-vector-ff-width-dp91}
\end{equation}
因此 $W$、$Z$ 的轻费米子宽度只需对同一个一般结果代入不同的手征系数。

对
\begin{equation*}
W^-(P,\lambda)\to\ell^-(p)+\bar\nu_\ell(k),
\end{equation*}
带电流只有左手分量，因此
\begin{equation}
 c_L^{W}=\frac{g}{\sqrt2},\qquad c_R^{W}=0,
 \qquad
 \mathcal M_\lambda
 =\frac{g}{\sqrt2}\epsilon_\mu^{(\lambda)}(P)
 \bar u(p)\gamma^\mu P_Lv(k).
 \label{eq:wlnu-amplitude-dp11}
\end{equation}
将它代入\eqnrefs{eq:generic-vector-ff-width-dp91}，每个轻轻子族得到
\begin{equation}
\boxed{
\Gamma(W\to\ell\nu)
=\frac{g^2M_W}{48\pi\hbar}
=\frac{G_F^{(E)}M_W^3}{6\pi\sqrt2\,\hbar}.}
\label{eq:w-leptonic-width-dp11}
\end{equation}
若末态是 $u_i\bar d_j$，只需再代入颜色多重度和 CKM 权重：
\begin{equation*}
\Gamma(W\to u_i\bar d_j)
=3|V_{ij}|^2\Gamma(W\to\ell\nu).
\end{equation*}
在足够远离重夸克阈值的区域，微扰量子色动力学再给出 $1+\alpha_s(M_W)/\pi+\cdots$ 的短程修正。

对中性流
\begin{equation*}
Z(P,\lambda)\to f(p)+\bar f(k),
\end{equation*}
正文的矢量--轴矢量顶角
\begin{equation}
\mathcal M_\lambda
=\frac{g}{2c_W}\epsilon_\mu^{(\lambda)}(P)
\bar u(p)\gamma^\mu(g_V^f-g_A^f\gamma^5)v(k)
\label{eq:zff-amplitude-width-dp11}
\end{equation}
可直接改写成手征形式
\begin{equation*}
 c_L^Z=\frac{g}{2c_W}(g_V^f+g_A^f),\qquad
 c_R^Z=\frac{g}{2c_W}(g_V^f-g_A^f).
\end{equation*}
于是\eqnrefs{eq:generic-vector-ff-width-dp91} 给出
\begin{equation}
\boxed{
\Gamma(Z\to f\bar f)
=N_c^f\frac{G_F^{(E)}M_Z^3}{6\pi\sqrt2\,\hbar}
\left[(g_V^f)^2+(g_A^f)^2\right].}
\label{eq:zff-width-tree-dp11}
\end{equation}
这里 $N_c^f=1$ 对轻子，$N_c^f=3$ 对夸克。靠近重费米子阈值时，必须恢复末态质量并重新计算相空间与手征干涉；\eqnrefs{eq:generic-vector-ff-width-dp91} 只对应 $m_f/M_V\to0$ 的受控极限。

\begin{figure}[htbp]
\centering
\begin{tikzpicture}[line width=.8pt,line label/.style={fill=white,fill opacity=.96,text opacity=1,inner sep=1.15pt}]
\begin{scope}[xshift=-3.2cm]
\draw[decorate,decoration={snake,amplitude=1.4pt,segment length=6pt}] (-2,0)--(0,0) node[line label,midway,above]{$W^-$};
\draw[fermion] (0,0)--(2,1) node[line label,midway,above]{$\ell^-$};
\draw[fermion] (2,-1)--(0,0) node[line label,midway,below]{$\bar\nu_\ell$};
\end{scope}
\begin{scope}[xshift=3.4cm]
\draw[decorate,decoration={snake,amplitude=1.4pt,segment length=6pt}] (-2,0)--(0,0) node[line label,midway,above]{$Z$};
\draw[fermion] (0,0)--(2,1) node[line label,midway,above]{$f$};
\draw[fermion] (2,-1)--(0,0) node[line label,midway,below]{$\bar f$};
\end{scope}
\end{tikzpicture}
\caption{$W$ 与 $Z$ 的树级费米子衰变。同一组电弱流顶角在低能虚交换中形成四费米子算符，在质量壳上则决定规范玻色子的部分宽度。}
\label{fig:wz-fermionic-decays-dp11}
\end{figure}
